\subsection{Vector tests for lines in space}

Let
\[
L_1=P_1+t\mathbf v_1,
\qquad
L_2=P_2+s\mathbf v_2.
\]
The directions are parallel precisely when
\[
\boxed{
\mathbf v_1\times\mathbf v_2=\mathbf0.
}
\]
This is the three-dimensional parallelism test established in the vector
chapter. It determines whether the directions agree, but it does not by itself
distinguish identical lines from distinct parallel lines. For that, one must
also check whether $P_2-P_1$ is parallel to the common direction.

The direction vectors are perpendicular precisely when
\[
\mathbf v_1\cdot\mathbf v_2=0.
\]
However, in space this alone does not imply that the geometric lines are
perpendicular. Two lines are perpendicular only when they both intersect and
their directions form a right angle. Without intersection they may be skew even
though their direction vectors are perpendicular.

\begin{examplebox}
Consider
\[
L_1=(0,0,0)+t(1,0,0)
\]
and
\[
L_2=(0,1,1)+s(0,1,0).
\]
Their direction vectors satisfy
\[
(1,0,0)\cdot(0,1,0)=0.
\]
Yet a point on $L_1$ has third coordinate $0$, while every point on $L_2$ has
third coordinate $1$. The lines cannot intersect. They have perpendicular
directions but are skew, not perpendicular lines.
\end{examplebox}

To decide whether two nonparallel lines intersect, solve
\[
P_1+t\mathbf v_1=P_2+s\mathbf v_2.
\]
This produces three scalar equations for the two parameters. A common pair
$(t,s)$ gives the intersection point; an inconsistent system proves that the
lines are skew.

\begin{center}
\begin{tikzpicture}[scale=0.9,>=stealth]
    \draw[blue!70!black,line width=1.1pt] (0.2,0.8) -- (5.4,0.8)
        node[right] {$L_1$};
    \draw[orange!85!black,line width=1.1pt] (2.0,2.0) -- (2.0,4.8)
        node[above] {$L_2$};
    \draw[gray,dashed] (2.0,2.0) -- (2.0,0.8);
    \node[gray,anchor=west] at (2.15,1.4) {separation in space};
    \draw[blue!70!black,->] (1.0,0.8) -- (2.2,0.8)
        node[midway,below] {$\mathbf v_1$};
    \draw[orange!85!black,->] (2.0,2.4) -- (2.0,3.6)
        node[midway,right] {$\mathbf v_2$};
\end{tikzpicture}
\end{center}
